% =======================================
% MEETING 2
% ---------------------------------------
% DATE   - AUGUST 14, 2026
% FORMAT - FACE TO FACE
% =======================================
% TOPIC  - EQUIVALENCE PRINCIPLE
% =======================================

\subsection{Einstein's Equivalence Principle}

	\subsubsection{Inertial Frames}

		The Weinberg definition is that, at every point in spacetime, i.e. at every $(t,x,y,z)$ 
		in a gravitational field, it is possible to choose a locally inertial frame, i.e. a freely 
		falling coordinate system, such that the ``laws of nature'' take the same form as they do in 
		an unaccelerated Cartesian coordinate system in the absence of gravity.

		In other words, fields that are accelerating are locally indistinguishable from gravity. 
		This is the basic idea behind the equivalence principle.

		Gravity can therefore be regarded as a fictitious force, because it can be removed by a 
		coordinate transformation. It is there only because we are using the wrong coordinate system.

		The equivalence principle states that it is always possible to remove gravity locally. 
		Therefore, gravity is a fictitious force'' that can be ``removed'' when describing the 
		local field in an experiment.

		The important difference between gravity and other fictitious forces is that gravity can 
		only be removed locally, while other fictitious forces can be removed globally. 
		The reason is that a gravitational field has spatial variation. Even after transforming 
		away the gravitational acceleration at one point, neighboring points may experience 
		slightly different accelerations.

		We say that, regardless of the mass of the object, it does not take part in determining 
		its own gravitational acceleration. The gravitational acceleration is universal.

		The measurement of differences between these gravitational fields is called \textit{tidal} measurement, 
		or simply the measurement of tides. These tidal effects cannot be transformed away over an extended 
		region.

	\subsubsection{Universality of Free Fall}

		The usual way this is expressed is to imagine a person in a closed box with no reference to the outside world. To gain 
		information about his environment, he begins conducting a free-fall experiment. Suppose they perform some experiment in 
		which they can measure the acceleration of objects.

		For an object $o_1$, they measure

		\begin{equation}
			\vb{a}_1 = \frac{m_{1g}}{m_{1i}} \vb{g}
		\end{equation}

		and for a second object $o_2$, they can measure

		\begin{equation}
			\vb{a}_2 = \frac{m_{2g}}{m_{2i}} \vb{g}.
		\end{equation}

		The universality of free fall is the statement that

		\begin{equation}
			\vb{a}_1 = \vb{a}_2
		\end{equation}

		under the same circumstances, regardless of the nature or mass of the objects. This is called the \textit{universality of free fall}.
		
		Suppose that the box is accelerating upward with acceleration $g\vu{z}$.
		\begin{equation}
			\vb{a}_\text{box} = g\vu{z}
		\end{equation}

		The observer releases the ball. From his perspective, the ball appears to accelerate downwards.
		\begin{equation}
			\vb{a}_\text{ball} = -g\vu{z}
		\end{equation}
		Therefore, it is reasonable to conclude that ``there is a gravitational field pulling the force downward".

		But there isn't. The ball is actually just continuing along its inertial trajectory while the box accelerates upward.

		Now compare this with another box which is stationary, with another observer sitting on the surface of Earth.

		The observer releases the same ball, and it falls downward with approximately

		\begin{equation}
			\vb{a}_\text{ball} = -g\vu{z}
		\end{equation}
		Inside the box, the experiment looks the same. So locally, motion due to acceleration is indistiguishable from motion
		incduced by a gravitational field. There is no experiment that can be performed by the two observers that can differentiate
		one from another. This is the strong equivalence principle.

		In Newtonian language, this corresponds to the equality of gravitational and inertial mass. The important point for the equivalence principle is that all sufficiently small freely falling objects respond to the gravitational field in the same way.

		Now, imagine the observer is in a box or elevator that is freely falling towards Earth.
		The box and everything inside it are accelerating downward together.
		Suppose the observer release a ball inside the box.
		The ball is falling toward Earth, but so is the box. Since both have essentially the same gravitational acceleration,
		\begin{equation}
			\vb{a}_\text{ball}\approx \vb{a}_\text{box}
		\end{equation}
		their relative acceleration is approximately zero:
		\begin{equation}
			\vb{a}_\text{relative} = \vb{a}_\text{ball} - \vb{a}_\text{box} \approx 0
		\end{equation}
		so the ball appears to float. The observer could throw it, and it would move approximately in a straight line at constant velocity relative to the observer.

		The observer would therefore reasonably conclude ``There is no gravity here. I'm in an inertial frame."

		But the observer actually is in Earth's gravitational field. In actual inertial motion, far from any gravitating bodies, the same observations will take place.
		Therefore, a falling observer is indistinguishable from an inertial observer. That is the weak equivalence principle.


		In a simple example, we can write the gravitational field as a non-inertial frame:

		\begin{equation}
			\ddot{y} = - g
		\end{equation}

		We can remove gravity, however, via certain transformations. Let

		\begin{equation}
			z = y - \frac{1}{2}gt^2
		\end{equation}

		This results in

		\begin{equation}
			\ddot{z} = 0.
		\end{equation}

		This is an inertial coordinate system, since there is no perceived acceleration. The point is that gravity can be removed locally by choosing an appropriate freely falling coordinate system.

		The person in the box cannot tell whether they are accelerating upward with constant acceleration $g$, or whether they are freely falling under the influence of gravity. In a sufficiently small region, the two situations are experimentally indistinguishable.

		Gravity is therefore, locally, simply the manifestation of being in the wrong frame.

		\textit{Inertial} means the absence of gravitational acceleration in the local frame. More precisely, an inertial frame is one in which freely falling objects move without coordinate acceleration.

		\paragraph{Weak, Strong, and Einstein Equivalence Principles}

			We have the weak equivalence principle, where the ``laws of nature'' are understood primarily as the laws of mechanics and the universality of free fall.
			For the strong equivalence principle, one extends the statement to gravitationally self-interacting or strongly gravitating systems, so that the equivalence applies more broadly than to test particles alone.

			The Einstein equivalence principle goes further: ``laws of nature'' genuinely means all laws of nature. It is one of the conceptual foundations of general relativity.
			\begin{theorem}{The Principle of Relativity}
				\label{thm:principle_of_relativity}
			    The laws of nature take the same form in all inertial observers.
			\end{theorem}
			
			In any theory of gravity, one can attempt to remove $g$ locally by an appropriate choice of coordinates. What makes general relativity special is the geometric interpretation of what remains after the local gravitational acceleration has been removed.

			It means that gravity is genuinely fictitious locally if no experiment confined to that sufficiently small region can tell that we are in a gravitational field. If one can remove gravity, then what remains should be the stage itself: the underlying spacetime structure and its geometry.

			In summary, over a sufficiently small region of space and time, a fictitious gravitational field produced by acceleration is indistinguishable from a gravitational field produced by mass (strong equivalence principle).
			Likewise, within a sufficiently small region of space, a freely falling observer cannot distinguish between free fall in a gravitational field and inertial motion in the absence of gravity (weak equivalence principle). 
			Therefore, a freely falling observer is locally an inertial observer.
\subsection{History of Relativity}

	The notion of inertiality plays a major role in relativity. Let us first explain relativity in its historical context.

	We begin with special relativity and Newtonian gravity. The combination of these ideas eventually leads us toward general relativity, where gravity is incorporated into the structure of spacetime itself.

	\subsubsection{Relativity before Einstein}

		What is relativity? It is an old notion that existed even by the time of Newton. The essential idea is that there should exist observers for which the laws of nature take the same form. Regardless of where the observer is or how they move within the appropriate class of frames, the laws of nature should be invariant.

		This is the equivalence of \textit{inertial observers}: the laws of nature take the same form for all inertial observers.

		Aristotle did not have this notion. He believed that there was a special place toward which objects naturally moved: the center of the universe, identified in practice with the center of the Earth. One could therefore imagine that this special region defined the natural frame in which the laws of nature took their proper form.

		It made sense within that picture to express physics relative to the center of the Earth, since all other frames could be considered as moving relative to this preferred location. For example, one could imagine waves in a pond as being fixed relative to the pond itself, with the pond providing a preferred rest frame.


		Galileo rejected this notion by positing that there is not a single preferred inertial state. Instead, there is a class of observers for whom the laws of mechanics are equally valid. These observers move uniformly relative to one another.

		This is called \textit{Galilean relativity}. There is a class of observers which all move relative to each other with uniform velocities.

		The transformation between two such frames can be written as
		\begin{align}
			t' &= t\\
			\vb{x}' &= \vb{x} - \vb{v}t
		\end{align}

		This relation explains how one can switch to the frame of reference of a second observer moving at velocity $\vb{v}$ relative to the first.

		In that case, one has the following invariance of Newtonian mechanics:

		\begin{equation}
			\vb{F} = m\vb{a} \implies \vb{F}' = m\vb{a}'.
		\end{equation}

		Newtonian theory therefore satisfies Galilean relativity: it is invariant under Galilean transformations.

		From the development of Newtonian mechanics onward, classical mechanics became the fundamental framework of physics.
		From 1687 onwards, the only fundamental laws of physics is Newtonian physics.

	\subsubsection{Maxwell's Results}

		For roughly three centuries, Newtonian mechanics was the jewel of physics. Then electromagnetism broke this picture.

		In 1862, Maxwell showed that Maxwell's equations are not Galilean invariant. They therefore do not satisfy Galilean relativity in the same way that Newton's equations do. This challenged centuries of orthodoxy.

		The response of academia was initially to deny that Maxwell's equations were fundamental in the same sense as Newtonian mechanics. One possibility was that Maxwell's equations were only correct in the rest frame of some physical medium, and that medium was called the \textit{aether}.

		The aether was imagined as an invented medium through which electromagnetic waves propagated. This was reasonable at the time, since
		all known waves needed a medium to travel from one place to another, think of sound or earthqakes or water waves.
		In this picture, Maxwell's equations described electromagnetic waves in the same broad conceptual sense that waves 
		in a pond are described relative to the pond.

		Light, therefore, would then only move at speed $c$ only in the rest frame of the aether.

		This led naturally to attempts to detect the Earth's motion through the aether, such as the Michelson-Morley experiment. However, the results of the experiment was null; the expected difference in the measured speed of light was not observed, and the result indicated that $c$ was the same in all different experimental directions.

		People fought against this result, as they should have. The natural response to an experimental discrepancy is to try to preserve a successful theory if possible and search for a physical explanation for the new result.

	\paragraph{Ad Hoc Solutions to the Aether Problem.}

		One proposed solution was Lorentz contraction. Objects would contract in the direction of motion through the aether, although there was initially no satisfactory physical explanation for why they should do so. It was essentially an ad hoc solution introduced to maintain consistency with the null result.

		Another proposed idea was \textit{ether dragging}. Imagine that the ether is not a fixed frame, but behaves something like molasses or syrup, so that the Earth in motion drags the ether around it. The Earth would then remain locally at rest with respect to the ether, and there would consequently be no observed change in $c$.

		The debates went on for decades. People had built careers around various versions of these theories, and many ad hoc solutions were introduced to preserve the compatibility of Maxwell's theory with Galilean relativity.

		There were dozens of theories introduced in the hope of making Maxwell's equations consistent with the Galilean framework.

	\subsubsection{Einstein's Contribution}

		Einstein came in 1905 with his \textit{annus mirabilis} papers.

		The three major papers associated with this period gave rise to major developments in quantum theory, statistical mechanics, and special relativity. Einstein developed these ideas all before turning 24.

		Everyone had been focused on finding a way to mesh Galilean relativity with Maxwell's equations. The common strategy was to deny the fundamentality of Maxwell's equations by claiming that they only worked in the preferred frame of the aether.

		Einstein instead reversed the argument.

		Einstein posited that Maxwell's equations are fundamental, while Newtonian mechanics is not fundamental. Newtonian mechanics was a theory that had been successful for three centuries, but its enormous historical success did not make it fundamentally immune to revision.

		His solution is:

		\begin{itemize}
			\item adopt $c$ as universal;
			\item accept Maxwell's equations as fundamental.
		\end{itemize}

		As such, the equations governing Galilean relativity must be dropped so that $c$ remains constant. By consequence, one also has to give up or modify Newtonian mechanics and stop treating it as fundamentally exact.

		One has to divorce oneself from the notion that Newtonian mechanics is fundamental.

		The notion of universal observers is still true: there is a class of observers for whom the laws of motion retain the same form. However, to reconcile this with the invariance of $c$, one must change the transformation law between observers.

		This leads to the unification of space and time and to the Lorentz transformations.

	\paragraph{Lorentz Transformations and Covariance}

		Maxwell's theory is relativistic, meaning that is relativistically covariant: under the appropriate relativistic transformations, its mathematical form does not change.
		This covariance is not apparent in the separate electric and magnetic fields, which transform into one another between inertial frames. Relativity instead unifies them as aspects of a single electromagnetic field.

		The task is therefore to find the correct transformation laws that keep Maxwell's equations invariant.
		These transformations will replace the Galilean transformation laws and reveal the kinematic structure required by electromagnetism.

		It is then found that the Poincar\'e-Lorentz transformations, initially formulated to explain the null result of Michelson-Morley and to save the aether hypothesis, provide the appropriate transformations:
		\begin{theorem}{Lorentz Transformations}
			Transformations of coordinates that allow $c$ to be invariant across all inertial observers.
			\begin{align}
				ct' &= \gamma(ct-\beta x) \\
				x'  &= \gamma(x - \beta ct) \\
				\gamma &= \frac{1}{\sqrt{1-\beta^2}} \\
				\beta &= \frac{v}{c}.
			\end{align}
		\end{theorem}

		The project, according to Einstein, is to rewrite all of physics so that it is in ``relativistically covariant'' form: the equations 
		should be invariant in form from one inertial observer to another\footnote{This isn't too difficult back in 1905, since the only fundamental theories of physics at the time were Newton's and Maxwells.}.

		The form of these transformations is
		\begin{equation}
			x^{\mu'} = \tensor{\Lambda}{^\mu_\nu} x^\nu.
		\end{equation}
		where
		\begin{equation}
			x^\mu = (ct, x^1, x^2, x^3)
		\end{equation}
		and
		\begin{equation}
			\tensor{\Lambda}{^\mu_\nu} =
			\bmqty{
				\gamma & -\gamma \beta & 0 & 0\\
				-\gamma \beta & \gamma & 0 & 0 \\
				0 & 0 & 1 & 0 \\
				0 & 0 & 0 & 1
			}
		\end{equation}

		Now, the force equation becomes
		\begin{equation}
			\dv{\vb{p}}{t} = \vb{F}\implies \dv{p^\mu}{\tau} = F^\mu.
		\end{equation}
		The second equation (the four momentum and the four force) reduces to the first equation (the Newtonian vector form),  in the appropriate nonrelativistic, low-velocity limit.

	\paragraph{How About Gravity?}

		Gravity is simply

		\begin{equation}
			\lpc\Phi = 4\pi G \rho.
		\end{equation}

		This is not relativistically covariant, since it is written as a theory involving only space and does not incorporate space and time into a unified relativistic structure.

		One can try to modify it so that

		\begin{equation}
			\square \Phi = 4\pi G \tensor{T}{^\mu_\nu}.
		\end{equation}

		This is known as Nordstr\"om gravity. But this theory is incorrect as a theory of gravity because it does not predict the bending of light.

		However, Einstein was bothered by the equivalence principle. Why is there an equivalence of falling? Why does gravity have this peculiar property that all objects fall in the same way?

		This question leads to the deeper geometric interpretation of gravity.

	\paragraph{Minkowski Spacetime}

		It was Minkowski in 1907 who placed these transformations in their modern geometric form.

		One can make sense of special relativity by imagining space and time as a single structure: \textit{spacetime}. This spacetime possesses a metric, now known as the Minkowski metric,

		\begin{equation}
			\eta_{\mu\nu} =
			\bmqty{
				-1 & 0 & 0 & 0 \\
				\phantom{-}0 & 1 & 0 & 0 \\
				\phantom{-}0 & 0 & 1 & 0 \\
				\phantom{-}0 & 0 & 0 & 1
			}
		\end{equation}

		and the transformations that preserve the Minkowski metric are the correct relativistic transformations.

		The key conceptual step is that space and time are no longer treated as two completely separate structures. They form a single spacetime geometry, and inertial observers correspond to the natural geometric structure of this spacetime.

		\subsubsection{From Special to General Relativity}

		In 1915, Einstein formalized his theory of General Relativity, where he formalized the notion of spacetime and recognized that spacetime need not be described only by the Minkowski metric.

		The replacement is

		\begin{equation}
			\eta_{\mu\nu}\to g_{\mu\nu}.
		\end{equation}

		The Minkowski metric is therefore promoted from the fixed metric of special relativity to the more general spacetime metric $g_{\mu\nu}$, which may vary from place to place.

		Einstein expanded the notion of inertial observers to all observers in free fall within spacetime. Inertial observers are therefore no longer merely observers moving with uniform velocity relative to one another. They are also freely falling observers, following geodesics of the spacetime geometry.

		He got the basic idea relatively early, but obtaining the field equations took much longer.

		Relativity is very, very old, and the important thing Einstein did was expand it into a general theory in which the geometric structure itself becomes dynamical.

		Minkowski's relativity is special relativity, represented by $\eta$, called special because it is restricted to a class of observers moving uniformly relative to one another.

		The general form, represented by $g$, becomes general relativity, with gravity incorporated into the geometry of spacetime.
	
		\begin{question}{Wasn't Einstein a Machian?}
			There is a distinction here involving the idea of being \textit{Machian}. A Machian picture suggests that everything should be relational and that nothing should be absolute. Einstein was a Machian, but the theory that he founded is fundamentally not Machian.
			General relativity has a deeply relational character, although it is not simply Machian in every possible sense. The geometry of spacetime itself becomes part of the physical description rather than serving merely as an immutable background stage.
		\end{question}
	\paragraph{Relativity is more general than Light}

		% There is a notion that light is inherently tied to relativity. However, even if Maxwell's theory were somehow removed or replaced, the principle of relativity itself would remain.

		% Light has nothing fundamentally to do with the concept of relativity. Its role is historical: Maxwell's equations were what forced physicists to reconsider Galilean relativity and search for a new transformation law, but the principle of relativity is conceptually distinct from electromagnetism.

		% Relativity is a metaprinciple. Maxwell inspired the modern development of relativity, but Maxwell's theory is not the conceptual foundation of the principle itself.

		% Light is therefore not intrinsically tied to relativity. One can derive general relativity without making light the starting point, and obtain the same gravitational equations.

		% The next lecture goes on to deriving general relativity while making no fundamental use of light, yet arrives at the same equations.


		There is a common tendency to regard light as being intrinsically tied to
		relativity. Historically, however, light played a more specific role. Even if
		Maxwell's theory were somehow removed or replaced by another theory, the principle of
		relativity would remain.\footnote{Sir here goes on a rant about news sites 
		which sensationalize `slowing down light speed' experiments as `putting to question
		relativity'}

		Light has no fundamental role in the concept of relativity itself. Its
		importance is primarily historical: Maxwell's equations revealed that
		electromagnetic waves propagate at a universal speed that is incompatible with
		Galilean kinematics. This compelled physicists to reconsider the Galilean
		transformation and ultimately led to the Lorentz transformation. The principle
		of relativity, however, is conceptually distinct from electromagnetism.

		Relativity therefore is a metaprinciple. It is a more general principle concerning the relationship between
		physical laws and different observers. Maxwell's theory inspired and helped reveal the
		structure of spacetime, but it is not the conceptual foundation of the
		principle of relativity itself.

		Light is therefore not intrinsically tied to relativity. The relativistic
		structure of spacetime can be developed without taking electromagnetism or
		light as its starting point. Likewise, general relativity can be formulated
		without making light fundamental to its derivation, while still arriving at
		the same gravitational field equations.

		The next lecture will take precisely this approach: we will develop 
		relativity without making any mention of light and arrive at
		the same equations that describe gravity in Einstein's theory.
	\paragraph{The Role of $c$}

		% The only essential role of Maxwell's theory in this historical development is to establish $c$ as a universal constant that characterizes the structure of spacetime.

		% We can approach relativity with different values of the invariant speed. If we assign
		% \begin{equation}
		% 	c=\infty,
		% \end{equation}
		% we recover Galilean relativity.\footnote{This was in fact the case during the time of Galileo. It was a mainstream hypothesis that the speed of light was infinite}

		% If we assign
		% \begin{equation}
		% 	c=0,
		% \end{equation}
		% we obtain Carrollian relativity\footnote{This is in reference to Lewis Carrol's \textit{Alice in Wonderland}, likening the bizarre events that happen in such conditions to the plot of said work}.

		% And if we assign $c$ to be an arbitrary finite number between $0$ and $\infty$,
		% \begin{equation}
		% 	c\in(0,\infty),
		% \end{equation}
		% we obtain the relativistic structure associated with special relativity, and ultimately the framework in which general relativity is formulated. Do note that
		% $c$ need not be a specific value, but it is a constant parameter that can take on any fixed finite nonzero value. It acts as the \textit{speed limit} of the universe, the invariant
		% speed of the relativistic kinematic structure. The quantity $c$ is called the \textit{speed of light} 
		% because light is the first entity observed to travel at such speed, but later on we realized that gravitational waves and other massless
		% excitations propagate at such speed, while massive particles can approach it arbitrarily closely but cannot reach it. 
		% Light only gave $c$ the numerical designation we know today as \SI{299792458}{m/s}. 

		% We then develop a theory, without any mention of light, and see that a \textit{speed limit} will fundamentally arise
		% from first principles. We will tackle this in the next chapter.

		The essential role of Maxwell's theory in the historical development of
		relativity is to identify a universal invariant speed, $c$, that characterizes
		the structure of spacetime. However, the existence of such an invariant speed
		does not depend specifically on electromagnetism. We can construct relativistic
		kinematics with any finite, nonzero value of the invariant speed.

		Consider the limiting cases. If we take
		\[
			c \to \infty,
		\]
		we recover Galilean relativity.\footnote{This was in fact the case during the time of Galileo. It was a mainstream hypothesis that the speed of light was infinite}
		In this limit, the effects associated with the
		relativity of simultaneity vanish, and time becomes absolute.

		If instead we take
		\[
			c \to 0,
		\]
		we obtain Carrollian relativity,\footnote{This is in reference to Lewis Carrol's \textit{Alice in Wonderland}. It references the Red Queen's race in \textit{Through the Looking-Glass}, where the Queen tells Alice, ``Now, here, you see, it takes all the running you can do, to keep in the same place." In Carrollian physics, because $c\to0$, lightcones collapse entirely, meaning particles cannot move through space no matter how much momentum they have.}
		the ultra-relativistic counterpart of the
		Galilean limit. In this limit, the causal structure becomes degenerate in such
		a way that spatial motion relative to an inertial observer is effectively
		frozen.

		If instead we asign $c$ to be any fixed finite value,
		\[
			c \in (0,\infty),
		\]
		we obtain the Lorentzian kinematic structure underlying special relativity
		and, ultimately, general relativity. The particular numerical value of $c$
		is therefore not essential to the structure of the theory; any fixed finite
		nonzero value could serve as the invariant speed, with the numerical value
		depending on the choice of units and on the physical theory being described.

		The constant $c$ is commonly called the \textit{speed of light} because electromagnetic
		radiation propagates at this invariant speed. Historically, the value of $c$
		emerged from Maxwell's theory as the propagation speed of electromagnetic
		waves. Modern physics, however, understands $c$ more fundamentally as a
		property of spacetime itself: light travels at $c$ because photons are
		massless, just as gravitational waves propagate at $c$. Massive particles
		cannot reach $c$, although they may approach it arbitrarily closely.

		Thus, Maxwell's theory is not required in order to derive the existence of a
		finite invariant speed. Rather, it was the first physical theory to reveal
		that such a speed exists in nature and to give it the numerical value.
		We can therefore develop a theory of relativity without making any reference
		to light at all. Starting from the principles of spacetime symmetry, we will
		see that a finite invariant speed can arise from the kinematic structure
		itself. We will develop this approach in the next chapter.















